\section{Deformation coordinates and specialization}
\label{sec:deformations}

At a characteristic-$p$ point of formal height $1\leq h<n$, the full
height-$n$ $p$-divisible group has an \'etale Tate module of rank $n-h$.
We embed this module in the height-$n$ Honda universal cover and use a
specialization map to detect whether its vectors remain in the Tate
lattice at a nearby parameter point.  The values of this map on a Tate
basis will give analytic functions with invertible differential along
the height-at-least-$h$ parameter space.

We first construct the specialization map and prove this differential
statement.  An integral comparison with the height-$h$ Honda group then
gives the nonzero natural map needed in
Proposition~\ref{bundles:nondegeneracy}.  Finally, we determine how the
specialization map changes under the stabilizer action.

\subsection{The universal family and its height strata}

Let $H$ be the universal deformation of $\Gamma_n$ over
\[
 E_0=W(k_n)[[u_1,\ldots,u_{n-1}]],
 \qquad
 R_0=E_0/p=k_n[[u_1,\ldots,u_{n-1}]].
\]
We also write $H$ for its reduction over $R_0$.
Choose the usual successive height parameters.  The Lubin--Tate parameter
construction and universality theorem
\cite[Proposition~1.1 and Theorem~3.1]{LT66} give the following convention:
after setting $p,u_1,\ldots,u_{h-1}$ to zero, the $p$-series factors as
\begin{equation}
 [p]_H(T)=f(T^{q_h}),\qquad
 f(0)=0,\qquad
 f'(0)\in u_h(E_0/I_h)^\times
 \quad (h<n),
 \label{def:pseries}
\end{equation}
where $f$ is considered over $E_0/I_h$ and the coefficient is a unit when
$h=n$.  Equivalently, the successive leading coefficients of the
$p$-series generate the ideals
$I_h=(p,u_1,\ldots,u_{h-1})$.

We recall why the factorization assertion holds over nonreduced
characteristic-$p$ rings as well.  For a homomorphism $a:F\to F'$ of formal
groups, differentiation of the homomorphism identity gives
\[
 a'(T)\,\partial_2F(T,0)
   =\partial_2F'(a(T),0)\,a'(0).
\]
Both displayed partial derivatives are units.  Thus $a'(0)=0$ implies
$a'(T)=0$, and $a(T)$ factors through $T^p$.  The resulting series is a
homomorphism from the Frobenius twist of $F$.  Applying this argument
successively to the $p$-series proves the factorization in
Equation~\eqref{def:pseries}.  It also shows that the condition defining $I_h$
is invariant under change of coordinate.  The universal deformation
isomorphisms therefore preserve these ideals.

Fix for the rest of this section
\[
 1\le h<n,\qquad r=n-h,\qquad q=p^h,\qquad Q=p^n.
\]
The height-at-least-$h$ parameter polydisc has coordinates
$y=(u_h,\ldots,u_{n-1})$, all of norm less than one; the preceding parameters
are zero.  Write $H_y$ for the evaluated law and
\[
 P_y(T)=[p]_{H_y}(T)=f_y(T^q),\qquad b(y)=f_y'(0).
\]
All their coefficients are integral and analytic in the parameters.
Choose a point $x$ in this polydisc with $u_h(x)\ne0$.  The formal group
$H_x$ over $C$ has height $h$, and $b(x)\ne0$.

\begin{lemma}\label{def:algebraization}
The finite group schemes
\[
 \mathcal G[p^j]
   =\Spec\bigl(R_0[[T]]/([p^j]_H(T))\bigr),\qquad j\ge1,
\]
form an algebraic $p$-divisible group $\mathcal G$ of height $n$ over $R_0$.
At the point $x$, its connected formal group is $H_x$ and its \'etale
height is $r=n-h$.  All the $C$-valued points of every level have coordinate
in $\mm_C$.
\end{lemma}

\begin{proof}
The preparation used here follows from a short division argument.
Let $(A,\mathfrak a)$ be a complete Noetherian local ring, let $d\geq1$, and let
$f(T)=T^d+e(T)$ with $e\in\mathfrak a A[[T]]$.
Let $\tau$ discard terms of degree below $d$ and shift the remaining
degrees down by $d$.  Then
\[
 \tau(qf)=(1+L)q,\qquad L(q)=\tau(qe).
\]
The operator $L$ raises $\mathfrak a$-adic order, so
$(1+L)^{-1}=\sum_{i\geq0}(-L)^i$ converges on $A[[T]]$.
For any $g$, take $q=(1+L)^{-1}\tau(g)$; the unique remainder
$g-qf$ has degree less than $d$.  Dividing $T^d$ gives
$T^d=qf+r$ with $q\equiv1\pmod{\mathfrak a}$ and
$r\in\mathfrak a A[T]$ of degree less than $d$.
Thus $q$ is a unit and $f=q^{-1}(T^d-r)$ is prepared.
The same division proves that $A[[T]]/(f)$ is free over $A$ on
$1,T,\ldots,T^{d-1}$.

Let $\mathfrak m=(u_1,\ldots,u_{n-1})\subset R_0$.  Reduction modulo
$\mathfrak m$ sends $[p^j]_H(T)$ to $T^{Q^j}$.  Weierstrass preparation
therefore writes
\begin{equation}
 [p^j]_H(T)=U_j(T)W_j(T),
 \label{def:prepared-level}
\end{equation}
with $U_j\in R_0[[T]]^\times$ and $W_j$ a distinguished monic polynomial
of degree $Q^j$.  The level algebra is consequently finite free of rank
$Q^j$.  It is $\mathfrak m$-adically complete, and its coordinate $T$ is
topologically nilpotent.  Substitution of the formal group law in the
finite level algebras thus defines genuine group-scheme operations.

Apply preparation once more to $[p]_H(T)-V$ over $R_0[[V]]$.  This shows
that substitution $V=[p]_H(T)$ is finite free of rank $Q$.  Reducing by
$[p^j]_H(V)$ proves that the multiplication maps between consecutive
levels are finite faithfully flat of that rank.  The kernel and inclusion
identities follow from the identities of the $p$-series.  These are exactly
the $p$-divisible-group conditions.

After evaluation at $x$, all nonleading coefficients of $W_j$ have norm
less than one.  At a point of norm at least one its leading term strictly
dominates the others, so every root is small.  The evaluated $U_j$ is a unit
on the open unit disc: its constant coefficient has norm one and its
remaining coefficients have norm at most one.  The roots of $W_j$ are
therefore precisely the small zeros of $[p^j]_{H_x}$, with the same group
operations.

The initial degree of $[p^j]_{H_x}$ is $q^j$, with nonzero coefficient.
Completing the finite algebra at the identity thus gives
$C[[T]]/(T^{q^j})$, with the group law induced by $H_x$.  These kernels
recover the formal group $H_x$: their defining ideals are cofinal with
the powers of $(T)$.  Its height is $h$, so the \'etale height of
the full height-$n$ group is $n-h$.
\end{proof}

\subsection{A natural specialization map}

Recall that
$\mathcal H_n(R)=\Gamma_n(R^{\circ\circ})$ for an affinoid perfectoid
$C$-algebra $R$.  We use an equivalent complete spectral norm on $R$;
it is power-multiplicative.  In particular $R^{\circ\circ}$ is its open
unit ball.  The estimates below use power-multiplicativity, not
multiplicativity on products of distinct elements.

\begin{proposition}[Specialization]\label{def:transport}
For each parameter $y$ in the height-at-least-$h$ polydisc, the formula
\begin{equation}
 \Theta_y(z)=\lim_{j\to\infty}
       P_y^{\circ j}\bigl(z^{1/Q^j}\bigr)
 \label{def:theta-limit}
\end{equation}
defines a continuous natural homomorphism
\[
 \Theta_y:\mathcal H_n(R)\longrightarrow H_y(R^{\circ\circ}).
\]
Here the roots are unique because $R$ is perfect, and $P_y^{\circ j}$
denotes composition.  The map
\begin{equation}
 \Phi_y:z\longmapsto
       \bigl(\Theta_y(p^{-m}z)\bigr)_{m\ge0}
 \label{def:cover-identification}
\end{equation}
is a natural isomorphism
\[
 \mathcal H_n(R)\xrightarrow{\ \sim\ }
     \varprojlim_{[p]_{H_y}} H_y(R^{\circ\circ}),
\]
and $\Theta_y$ is its zeroth projection.  For each fixed
$z\in\mathcal H_n(C)$, the function $y\mapsto\Theta_y(z)$ is analytic
on the parameter polydisc.
\end{proposition}

\begin{proof}
First fix $y$ and choose $\delta<1$ bounding the norms of its parameters.
The coefficients of $H_y-\Gamma_n$ and $P_y(T)-T^Q$ then have norm at most
$\delta$.  Integral power series are $1$-Lipschitz on the open unit ball.
Since $P_y(T)=f_y(T^q)$, power-multiplicativity gives
\begin{equation}
 \|P_y^{\circ j}(a)-P_y^{\circ j}(b)\|
       \le \|a-b\|^{q^j}
       \qquad(a,b\in R^{\circ\circ}).
 \label{def:contraction}
\end{equation}
For $m\ge0$ put
\[
 c_{j,m}=P_y^{\circ j}\bigl(z^{1/Q^{j+m}}\bigr).
\]
The discrepancy before the last $j$ compositions is bounded by $\delta$,
so Equation~\eqref{def:contraction} gives
\[
 \|c_{j+1,m}-c_{j,m}\|\le\delta^{q^j}.
\]
Thus $c_{j,m}$ converges, uniformly in $z$, and its limit satisfies
\begin{equation}
 \|\Phi_y(z)_m-z^{1/Q^m}\|\le\delta,
 \qquad
 \|\Phi_y(z)_m\|
       \le\max\{\|z\|^{1/Q^m},\delta\}<1.
 \label{def:cover-tail}
\end{equation}
Shifting the index in the limit proves
$P_y(\Phi_y(z)_{m+1})=\Phi_y(z)_m$.

Here is an explicit inverse.  If $(t_m)_{m\ge0}$ is a compatible sequence
in the inverse limit, define
\[
 \Psi_y((t_m))=\lim_{m\to\infty}t_m^{Q^m}.
\]
Indeed,
\[
 \|t_{m+1}^{Q^{m+1}}-t_m^{Q^m}\|
   =\|(t_{m+1}^Q-P_y(t_{m+1}))^{Q^m}\|
   \le\delta^{Q^m}.
\]
The limit $z$ therefore exists and obeys
\begin{equation}
 \|z-t_m^{Q^m}\|\le\delta^{Q^m},
 \qquad \|z\|\le\max\{\|t_0\|,\delta\}<1.
 \label{def:inverse-tail}
\end{equation}
No common bound less than one on the norms of all the $t_m$ is needed.
Equation~\eqref{def:cover-tail}, raised to the $Q^m$-th power, proves
$\Psi_y\Phi_y(z)=z$.  Conversely, Equation~\eqref{def:inverse-tail} at
$m+j$ gives
\[
 \|z^{1/Q^{m+j}}-t_{m+j}\|\le\delta.
\]
Applying Equation~\eqref{def:contraction} and using
$P_y^{\circ j}(t_{m+j})=t_m$ proves
$\Phi_y\Psi_y((t_m))=(t_m)$.

To verify the group law, the coefficients of $\Gamma_n$ lie in $k_n$, so
$Q$-th powers and their inverses commute with its addition.  For
$a=z^{1/Q^{j+m}}$ and $b=w^{1/Q^{j+m}}$, the two inputs
$\Gamma_n(a,b)$ and $H_y(a,b)$ differ by at most $\delta$.
After $j$ compositions their difference is at most $\delta^{q^j}$.
Since $P_y^{\circ j}$ is an $H_y$-homomorphism, passage to the limit gives
\[
 \Phi_y(\Gamma_n(z,w))_m
     =H_y(\Phi_y(z)_m,\Phi_y(w)_m).
\]
This proves additivity, and hence additivity of the inverse as well.
The uniform tail estimates and continuity of roots and convergent series
prove continuity of the maps.  All constructions commute with continuous
morphisms of perfectoid $C$-algebras.  They therefore give the asserted
natural maps, and commute with $\Zp$-multiplication by continuity from
integer multiplication.

Finally, fix $z\in\mathcal H_n(C)$ and a closed parameter polydisc
strictly inside the open polydisc.  Choose $\delta<1$ bounding all its
parameter norms.  Each finite stage of Equation~\eqref{def:theta-limit}
belongs to its Tate algebra: the substituted constant has norm less than
one and all coefficients of the law are bounded.  The same bound
$\delta^{q^j}$ holds in the parameter Gauss norm.  The stages consequently
converge in that complete Tate algebra.  This proves analyticity.
Partial differentiation is continuous on a fixed closed polydisc, with
operator norm bounded by the inverse coordinate radius.  The convergence
therefore also permits first differentiation.
\end{proof}

\begin{lemma}\label{def:tate-kernel}
At the chosen point $x$, the map $\Theta_x$ is surjective on $C$-points,
and
\[
 T_x=\ker\bigl(\Theta_x:\mathcal H_n(C)\to H_x(\mm_C)\bigr)
       \cong\Zp^r.
\]
In particular, a $\Zp$-basis $z_1,\ldots,z_r$ of $T_x$ is
$\Qp$-linearly independent in $\mathcal H_n(C)$.
\end{lemma}

\begin{proof}
Under Equation~\eqref{def:cover-identification}, the kernel consists of
sequences with $t_0=0$ and $[p]t_{j+1}=t_j$.  These are precisely the
compatible $p^j$-torsion points of $\mathcal G_x$.  By
Lemma~\ref{def:algebraization} every such point is small, and the
\'etale height is $r$.  Over the algebraically closed field $C$ the
resulting Tate module is $\Zp^r$.

For surjectivity, apply Weierstrass preparation to $[p]_H(T)-V$ over
$R_0[[V]]$.  After setting $y=x$ and $V=w\in\mm_C$, its distinguished
polynomial is monic with all lower coefficients small.  Algebraic
closedness gives a root, and every root is small.  Thus every small point
admits division by $p$ within the open unit disc.  Repeated choices give a
compatible sequence with any prescribed zeroth point, and
Proposition~\ref{def:transport} proves surjectivity of $\Theta_x$.
Finally, multiplying a putative rational relation among a $\Zp$-basis by
a sufficiently large power of $p$ proves the last assertion.
\end{proof}

\subsection{The parameter differential}

The values of $\Theta_y$ on a basis of $T_x$ will measure parameter
motion.  To show that their differential detects every direction, we
first identify parameter directions with marked first-order deformations
of the full $p$-divisible group.  This comparison uses the algebraic
$p$-divisible group over $R_0$, so it applies even though the map
$R_0\to C$ does not send the parameters to nilpotent elements.
For a $p$-divisible group $G$, write $G^\vee$ for its Cartier dual and
$\omega_G$ for its module of invariant differentials.

\begin{proposition}\label{def:ks}
The first-order parameter deformation of $\mathcal G$ induces an
isomorphism of free $R_0$-modules
\[
 \mathsf K:R_0^{n-1}\xrightarrow{\ \sim\ }
       \Hom_{R_0}(\omega_{\mathcal G^\vee},\Lie\mathcal G).
\]
For every small $C$-valued parameter point $x$, its specialization is the
linear isomorphism from parameter directions to marked first-order
deformations of $\mathcal G_x$, with the constant deformation as origin.
\end{proposition}

\begin{proof}
We use the square-zero deformation theorem of
\cite[Theorem~5.1]{Lau08}, in the cotangent-theoretic tradition of
Grothendieck and Illusie
\cite[Theorem~4.4 and Corollary~4.7]{Illusie85}, fixing its torsor structure
throughout.  If $B\to A$ has square-zero
kernel $J$ and $p$ is nilpotent, marked lifts of a $p$-divisible group $G/A$
form a torsor under
\[
 \Hom_A(\omega_{G^\vee},J\otimes_A\Lie G).
\]
The associated Kodaira--Spencer map is characterized by the
arbitrary-ring-lift equivariance of \cite[Equation~(5.1)]{Lau08}.
We use this same construction at the closed point and at every later
specialization.  The theorem applies in characteristic $p$ at every
prime, including $p=2$.

Apply this comparison to the two maps into
\[
 R_0[\epsilon_1,\ldots,\epsilon_{n-1}]/
       (\epsilon_i\epsilon_j:1\le i,j\le n-1)
\]
given by $u_i\mapsto u_i$ and $u_i\mapsto u_i+\epsilon_i$.
They define the $R_0$-linear map $\mathsf K$.
The source is the full module of parameter derivations: every series in
$R_0$ has a finite expression
\[
 a=\sum_{0\le e_i<p}u^e a_e^p,
\]
so ordinary differentials over $\Fp$ are free on $du_1,\ldots,du_{n-1}$.
The target is free of the same rank because $\Lie\mathcal G$ has rank one
and $\omega_{\mathcal G^\vee}$ has rank $n-1$.

At the closed point, it is enough to show injectivity of the tangent
map from formal-group deformations to $p$-divisible-group deformations:
both tangent spaces have dimension $n-1$, by
\cite[Proposition~2.6 and Theorem~3.1]{LT66} and the chosen square-zero
torsor theorem.  Suppose a parameter tangent vector has zero image.
Then its $p$-divisible deformation is marked-isomorphic to the constant
one, so its finite kernels have compatible marked isomorphisms.

Those kernels recover the formal-group isomorphism.  More explicitly,
for a deformation $F$ over a local Artin characteristic-$p$ algebra $B$
with maximal ideal $J$ and $J^a=0$, one has
\[
 ([p^j]_F(T))\subseteq(T^{p^j}),\qquad
 T^{aQ^j}\in([p^j]_F(T)).
\]
The first inclusion follows by iterated Frobenius factorization.
For the second, preparation gives a monic polynomial of degree $Q^j$
whose lower coefficients lie in $J$; in its quotient $T^{Q^j}$ lies
in $J$, so its $a$th power vanishes.  The kernel ideals are therefore
cofinal with the coordinate-adic topology.  Taking inverse limits of
the compatible finite Hopf-algebra isomorphisms yields a marked
formal-group isomorphism.  This is the formal reconstruction underlying
Tate's equivalence \cite[Section~2.2, Proposition~1 and Lemma~0]{Tate67}.
Lubin--Tate tangent injectivity forces the original vector to be zero.
Thus $\mathsf K$ is invertible modulo the maximal ideal of $R_0$.
Its determinant is a unit, so it is an isomorphism over $R_0$.

For a vector $v=(v_i)\in C^{n-1}$, the analytic first-order substitution
is the actual ring homomorphism
\[
 a\longmapsto a(x)+\epsilon\sum_i v_i\partial_i a(x)
       \quad\text{from }R_0\text{ to }C[\epsilon]/\epsilon^2.
\]
All evaluations converge because the coordinates of $x$ are small.
The arbitrary-map functoriality in the cited deformation theorem
identifies its deformation class with the specialization of $\mathsf K$.
This proves the final assertion.  In cotangent notation the specialized
map involves $\Omega^1_{R_0/\Fp}\otimes_{R_0}C$, not the intrinsic module
$\Omega^1_{C/\Fp}$, which is zero because $C$ is perfect.
\end{proof}

\begin{lemma}\label{def:height-rigidity}
Let $A$ be a local Artin $C$-algebra with residue field $C$, and let $F/A$
be a marked deformation of a height-$h$ formal group over $C$.
If $[p]_F(T)$ factors through $T^{p^h}$ over $A$, then $F$ is the constant
deformation, with its marking.
\end{lemma}

\begin{proof}
Factorization through $T^{p^h}$ is preserved by a coordinate change:
the $p^h$-th power of any series has only exponents divisible by $p^h$.
Choose an isomorphism of the special fiber with $\Gamma_h$ and apply the
marked Lubin--Tate classification over $A$.  It gives parameters
$a_1,\ldots,a_{h-1}$ in the maximal ideal of $A$.
Successively, modulo preceding parameters the coefficient of $T^{p^i}$
in the $p$-series is $a_i$ times a unit.  This is the height-parameter
property recalled above; it also follows directly from the leading law
term $a_i C_{p^i}$ in \cite[Proposition~1.1]{LT66}.  Here $C_{p^i}$ is the
reduction of the integral polynomial
$((X+Y)^{p^i}-X^{p^i}-Y^{p^i})/p$; its $p$-fold sum has coefficient
$p^{p^i-1}-1$, a unit modulo $p$.
The assumed factorization therefore forces every $a_i$ to vanish.
The zero parameter tuple is the constant deformation, and the
isomorphism supplied by \cite[Theorem~3.1]{LT66} respects the marking.
For $h=1$ there are no parameters and the same conclusion holds.
\end{proof}

\begin{proposition}[Tate coordinates]\label{def:coordinates}
For a $\Zp$-basis $z_1,\ldots,z_r$ of $T_x$, the analytic map
\begin{equation}
 \theta(y)=\bigl(\Theta_y(z_1),\ldots,\Theta_y(z_r)\bigr)
 \label{def:theta}
\end{equation}
has invertible differential at $x$ on the height-at-least-$h$ polydisc.
\end{proposition}

\begin{proof}
Analyticity follows from Proposition~\ref{def:transport}, and
$\theta(x)=0$.  Let $v$ be a direction in the kernel of its differential.
Write $A=C[\epsilon]/\epsilon^2$ and $y'$ for the corresponding
first-order parameter point.  Its preceding parameters remain zero.
By Proposition~\ref{def:ks}, it suffices to prove that
$\mathcal G_{y'}$ is the constant deformation of $\mathcal G_x$ with
its residue marking.  We will use the vanishing differential to lift the
Tate basis as compatible torsion sections, splitting the \'etale quotient.
The fixed-height condition will then trivialize the connected factor.

\smallskip
\noindent\emph{Lifting the Tate basis.}
For $j\ge1$ set
\[
 t_{i,j}(y)=\Theta_y(p^{-j}z_i),\qquad
 t'_{i,j}=t_{i,j}(y')\in A,
\]
where evaluation at $y'$ means first-order Taylor evaluation.

We verify membership in the finite group schemes before passing to a
$p$-divisible group.  Evaluate Equation~\eqref{def:prepared-level} at the
analytic function $t_{i,j}(y)$.  Its value at $x$ is small, so on a
neighborhood of $x$ this gives an identity of convergent analytic
functions
\[
 U_{j,y}(t_{i,j}(y))\,W_{j,y}(t_{i,j}(y))
       =\Theta_y(z_i).
\]
The first factor has nonzero reduced value at $x$.  Taylor evaluation in
$A$ makes the right side zero, by the choice of $v$, and the first factor
remains a unit.  Hence $W_{j,y'}(t'_{i,j})=0$: the Taylor coordinate is an
actual point of $\mathcal G_{y'}[p^j]$.
The analytic multiplication identities also give
\[
 [p]t'_{i,j+1}=t'_{i,j},\qquad [p^j]t'_{i,j}=0.
\]
These identities agree with the algebraic finite-level operations.
For any finite collection of Taylor points $a+\epsilon b$, equip $A$
with the complete submultiplicative norm
\[
 \|a+\epsilon b\|_s=\max\{|a|,s|b|\},
\]
choosing $s>0$ so that the points and parameter values are small.
The universal congruences defining the polynomial finite-level operations
then converge; their terms divisible by the appropriate source-level
polynomials vanish, so their evaluations agree with the Taylor-evaluated
formal operations.

Each finite identity is evaluated on a neighborhood where its finitely
many series converge.  Those neighborhoods may depend on $j$; all their
Taylor values lie in the same algebraic ring $A$.  In particular no bound
on the derivatives uniform in $j$ is required.

Sending a basis vector of $(\mathbb Z/p^j)^r$ to $t'_{i,j}$ gives a homomorphism
of finite group schemes.  These maps respect multiplication by $p$ and
also the level inclusions: the source inclusion sends the basis vector
at level $j$ to $p$ times that at level $j+1$, exactly as in the displayed
identity.  Over the strictly henselian local Artin ring $A$, the
connected--\'etale sequence exists and its \'etale quotient is the
unique constant lift of the residue quotient.  The constructed maps
induce an isomorphism onto that quotient at each level, since they do so
over $C$.  They therefore split it and give
\[
 \mathcal G_{y'}\cong
       \mathcal G_{y'}^0\times(\Qp/\Zp)^r.
\]

\smallskip
\noindent\emph{Trivializing the connected factor.}
At level $j$, first base-change the finite algebra, obtaining
$B_j=A[T]/(W_{j,y'}(T))$.  The idempotents of its residue algebra lift
uniquely across $\epsilon$, and select the identity factor $B_j^0$.
This factor is a finite local Artin algebra, so it equals its completion
at $(\epsilon,T)$.  Equation~\eqref{def:prepared-level} identifies that
completion with
\[
 B_j^0\cong A[[T]]/([p^j]_{H_{y'}}(T)).
\]
Indeed the preparation identity holds coefficientwise over $A$, and all
residue factors supported away from $T=0$ become units in this
completion.  The group law is $H_{y'}$, by Taylor evaluation of the same
preparation and group-law identities.  Since the preceding parameters
are zero and $b(x)\ne0$, the series $[p^j]_{H_{y'}}$ is $T^{q^j}$ times
a unit.  Thus the displayed algebra is also $A[[T]]/(T^{q^j})$.
This describes the connected factor by taking the identity factor after
algebraic base change.

Its formal group $H_{y'}$ satisfies the hypothesis of
Lemma~\ref{def:height-rigidity}, so it is the constant deformation of
$H_x$ with the residue marking fixed.  That formal isomorphism restricts
to all its finite kernels: the coordinates there are nilpotent, so every
formal series evaluates by a finite sum.  The connected $p$-divisible
factor is consequently constant with its marking.

\smallskip
\noindent\emph{Recovering the marked deformation.}
The splitting of the \'etale factor also has the required marking.
Over algebraically closed $C$, each connected finite group scheme has
only its identity as a $C$-point.  A map from a constant finite group
scheme is determined by the images of its $C$-points.  The actual
torsion points given by the chosen Tate basis therefore select the unique section of the
\'etale quotient on the special fiber.  The constructed splitting
reduces to this section.  Combining it with the marked connected
trivialization gives the constant deformation of the entire
$\mathcal G_x$, with the original marking.

Proposition~\ref{def:ks} now implies $v=0$.  Both the parameter tangent
space and the target of $D\theta_x$ have dimension $r=n-h$, so this
differential is invertible.
\end{proof}

\subsection{An integral comparison with the height-\texorpdfstring{$h$}{h} Honda group}

The next elementary recursion controls every coefficient of the forward
comparison.  Its necessary-and-sufficient form will also be used to
choose comparisons at individual orbit points.

\begin{lemma}\label{def:recursion}
Let $B$ be a characteristic-$p$ algebra, let
$f(Z)=bZ+c_2Z^2+\cdots\in B[[Z]]$, and put $q=p^h$.
For $\alpha(T)=\sum_{j\ge1}a_jT^j$, the identity
\[
 \alpha(f(T^q))=\alpha(T)^q
\]
holds if and only if, for every $j\ge1$,
\begin{equation}
 a_j^q-b^ja_j
       =\sum_{i<j}a_i[Z^j]f(Z)^i.
 \label{def:coefficient-recursion}
\end{equation}
If $B=C$, all coefficients of $f$ are integral, and the coefficients
$a_j$ satisfy these equations, then all $a_j$ are integral.
In particular, after prescribing any finite initial segment satisfying
the equations, every recursive extension by roots in $C$ has integral
coefficients.
\end{lemma}

\begin{proof}
Both sides of the identity contain only powers of $T$ divisible by $q$.
Their coefficients at $T^{qj}$ give exactly
Equation~\eqref{def:coefficient-recursion}, since the contribution of $a_j$
on the left is $b^ja_j$.  There are no other coefficient equations.

For the integrality assertion, argue inductively on $j$.  The right
side of Equation~\eqref{def:coefficient-recursion} is integral if the
preceding coefficients are integral.  A root with $|a_j|>1$ would make
$|a_j|^q$ strictly larger than both $|b^ja_j|$ and the norm of the right
side, which is impossible.  This includes $j=1$, whose right side is
zero.  The polynomial is monic, so roots exist in the algebraically
closed field $C$, and the same argument bounds every possible choice.
\end{proof}

\begin{proposition}\label{def:natural-comparison}
Choose a formal group isomorphism
$\alpha_x:H_x\xrightarrow{\sim}\Gamma_h$ over $C$.
Its forward coefficients are integral, and composition gives a nonzero
natural $\Qp$-linear morphism of sheaves on perfectoid spaces over $C$,
\begin{equation}
 \nu_x:\mathcal H_n\longrightarrow\mathcal H_h,
       \qquad \nu_x(z)=\alpha_x(\Theta_x(z)).
 \label{def:nu}
\end{equation}
\end{proposition}

\begin{proof}
The formal isomorphism exists by the height classification over the
algebraically closed field $C$.  It intertwines $[p]_{H_x}(T)=f_x(T^q)$
with $[p]_{\Gamma_h}(T)=T^q$.  Lemma~\ref{def:recursion} therefore shows
that every coefficient of $\alpha_x$ lies in $\OO_C$.

An integral series with zero constant term converges on
$R^{\circ\circ}$ for every affinoid perfectoid $C$-algebra $R$, with
value again in $R^{\circ\circ}$.  Its formal homomorphism identity
converges on those arguments as well.  Combined with
Proposition~\ref{def:transport}, this proves additivity and naturality of
$\nu_x$ on all such affinoids, hence as a sheaf morphism.
Continuity extends integer linearity to $\Zp$-linearity.  Multiplication
by $p$ is invertible on both Honda sheaves, so the morphism is
$\Qp$-linear.  If scalars are expressed as the sheaf of continuous
$\Qp$-valued functions, the same equality follows at geometric
valued-field points by naturality; these detect equality of coordinates
in a uniform affinoid algebra.  Equivalently,
Theorem~\ref{bundles:input} permits working throughout with additive
sheaf morphisms.

Write $a_1=\alpha_x'(0)\ne0$ and choose
$t\in C$ with $0<|t|<|a_1|$.  For $j\ge2$ the integral coefficients
satisfy
\[
 |a_jt^j|\le|t|^2<|a_1t|,
\]
so $\alpha_x(t)\ne0$.  Lemma~\ref{def:tate-kernel} supplies
$z\in\mathcal H_n(C)$ with $\Theta_x(z)=t$.
Thus $\nu_x(z)\ne0$, proving nonzeroness of the sheaf morphism.
Only the forward series has been evaluated on the whole open unit disc;
no convergence assertion about its inverse is needed.
\end{proof}

\subsection{Compatibility with the stabilizer action}

The coordinate construction is complete.  We now describe how the
specialization map and its Tate lattice change under the stabilizer.

The following finite-level statement concerns the full universal ring
for every $n\geq1$, independently of the chosen height stratum.

\begin{lemma}[Continuity of the stabilizer action]
\label{def:continuous-action}
Put $\mathfrak m_E=(p,u_1,\ldots,u_{n-1})\subset E_0$.
For $M\geq1$ and $N\geq M$, every element of $1+p^NA_n$ acts
trivially on $E_0/\mathfrak m_E^M$.  Consequently the stabilizer action
on $E_0$, and on its characteristic-$p$ quotient, is adically continuous.
\end{lemma}

\begin{proof}
Let $B=E_0/\mathfrak m_E^M$, let $J$ be its maximal ideal, and
write $F$ for the universal formal group over $B$.
For $d\in A_n$, lift its coefficients in $k_n$ to $B$, obtaining a
series $\widetilde d(T)$ with zero constant term.  Its homomorphism
defect
\[
 \delta(X,Y)=\widetilde d(F(X,Y))
       -_F F(\widetilde d(X),\widetilde d(Y))
\]
has all coefficients in $J$, because $d$ is a homomorphism on the
special fiber.  For every $a\geq1$, substitution in
$[p]_F(Z)=pZ+\sum_{j\geq2}c_jZ^j$ sends a series with coefficients
in $J^a$ and zero constant term to one with coefficients in
$J^{a+1}$.  This uses $p\in J$ and $2a\geq a+1$.
Thus $[p]_F^N(\delta)=0$ for $N\geq M$.

Since $[p]_F^N$ is itself a homomorphism, this identity says exactly
that $f_N=[p]_F^N\circ\widetilde d$ is an endomorphism of $F$.
Its special fiber is $p^Nd$, and
$\operatorname{id}_F-_F f_N$ has linear coefficient
$1-p^N\widetilde d'(0)$, a unit of $B$.
It is therefore an automorphism lifting $1-p^Nd$.
The marked Lubin--Tate classification now identifies the deformation
with its translate by $1-p^Nd$, so that this element fixes the
parameter map $E_0\to B$.
The assertion is uniform in $d$; replacing $d$ by $-d$ gives the
stated congruence subgroup.  These subgroups form a neighborhood basis
of the identity in $A_n^\times$, proving continuity.
\end{proof}

We use the stabilizer action convention
in which $g\in A_n^\times$ carries $x$ to a point $y_g$ together with an
integral isomorphism
\[
 \iota_g:H_x\xrightarrow{\ \sim\ }H_{y_g}
\]
whose closed-point reduction is $g$.
These are the universal deformation isomorphisms
\cite[Section~3.4]{LT66}.  The orbit remains in the height-at-least-$h$
polydisc.  Moreover, $y_g\to x$ as $g\to1$ by
Lemma~\ref{def:continuous-action}: a series in the $M$-th power of the
characteristic-$p$ parameter ideal evaluates with norm at most
$\|x\|^M$.  On the characteristic-zero generic fiber the same argument
uses $\rho^M$, where $\rho=\max\{|p|,\|x\|\}<1$.

\begin{lemma}\label{def:stabilizer}
For $g\in A_n^\times$ and $z\in\mathcal H_n(C)$ one has
\begin{equation}
 \Theta_{y_g}(gz)=\iota_g(\Theta_x(z)).
 \label{def:orbit-transport}
\end{equation}
Consequently $gT_x=T_{y_g}$.
\end{lemma}

\begin{proof}
The coefficients of $g$ lie in $k_n$, so
$g(z^{1/Q^j})=(gz)^{1/Q^j}$.  Choose $\delta<1$ bounding the evaluated
parameter ideals.  The coefficients of $\iota_g-g$ are then bounded by
$\delta$.  On small arguments the two series differ by at most that
amount.  By Equation~\eqref{def:contraction}, their images after $j$
compositions with $P_{y_g}$ differ by at most $\delta^{q^j}$.
Using $P_{y_g}^{\circ j}\iota_g=\iota_g P_x^{\circ j}$ and passing to
the limit proves Equation~\eqref{def:orbit-transport}.  It gives
$gT_x\subseteq T_{y_g}$; the inverse isomorphism gives equality.
\end{proof}

Choose a $\Zp$-basis $z_1,\ldots,z_r$ of $T_x$ as in
Lemma~\ref{def:tate-kernel}.  Since $r=n-h<n$,
Proposition~\ref{bundles:nondegeneracy} applies to this basis and $\nu_x$:
the subgroup
\[
 \bigl\{(\nu_x(dz_1),\ldots,\nu_x(dz_r)):d\in A_n\bigr\}
       \subset\Gamma_h^r(\mm_C)
\]
satisfies no nonzero formal power-series equation over $k$.
The invertible differential of Proposition~\ref{def:coordinates} and
the transport identity of Lemma~\ref{def:stabilizer} provide the other
inputs for the orbit argument in Section~\ref{sec:orbits}.
